\chapter*{คำนำ}
\addcontentsline{toc}{chapter}{คำนำ}

ตำรา \textbf{“ควอนตัมฟิสิกส์เบื้องต้น: การประยุกต์ภาพเสมือนจริงสามมิติและโครงข่ายประสาทเทียมตามหลักฟิสิกส์”} เล่มนี้ เรียบเรียงขึ้นเพื่อใช้ประกอบการเรียนการสอนรายวิชาฟิสิกส์ควอนตัมเบื้องต้น สำหรับนักศึกษาระดับปริญญาตรี คณะวิทยาศาสตร์และเทคโนโลยี มหาวิทยาลัยราชภัฏรำไพพรรณี ตลอดจนสถาบันอุดมศึกษาที่มีการจัดการเรียนการสอนด้านฟิสิกส์ ฟิสิกส์ประยุกต์ วิศวกรรมศาสตร์ และวิทยาการคอมพิวเตอร์

การก้าวเข้าสู่ยุคเทคโนโลยีควอนตัมรุ่นที่สอง (Quantum 2.0) ได้เปลี่ยนกระบวนทัศน์จากการศึกษาพฤติกรรมควอนตัมแบบสังเกตเฉย (Passive Observation) สู่การควบคุมและจัดการสถานะควอนตัมเชิงรุก (Active Manipulation) เช่น ระบบคิวบิต ความพัวพันเชิงควอนตัม และคอมพิวเตอร์ควอนตัม ผู้เขียนจึงได้ปรับปรุงและพัฒนาเนื้อหาในตำราเล่มนี้ให้มีความทันสมัย ครอบคลุมทั้งรากฐานทางทฤษฎีกลศาสตร์คลื่นดั้งเดิม (สมการชเรอดิงเงอร์ บ่อศักย์ ฮาร์มอนิกออสซิลเลเตอร์ อะตอมไฮโดรเจน) และสารสนเทศควอนตัมร่วมสมัย

จุดเด่นสำคัญของตำราเล่มนี้คือการบูรณาการนวัตกรรมการเรียนรู้เชิงรุก (Active Learning) 3 มิติ ได้แก่
\begin{enumerate}
    \item \textbf{ระบบภาพเสมือนจริงสามมิติและการติดตามภาพ (Augmented Reality \& Image Tracking):} ภาพประกอบในหนังสือได้รับการออกแบบให้เป็นเป้าหมายติดตามภาพ (Image Target) นักศึกษาสามารถส่องกล้องสมาร์ทโฟนเข้าหาภาพประกอบเพื่อเรนเดอร์โมเดล 3 มิติเชิงปฏิสัมพันธ์และแอนิเมชันคลื่นลอยตัวเหนือหน้าหนังสือได้ทันทีโดยไม่ต้องติดตั้งแอปพลิเคชัน
    \item \textbf{รหัสตอบสนองเร็วเพื่อเปิดห้องปฏิบัติการเสมือน (Vector QR Code Virtual Labs):} เชื่อมโยงสู่ WebXR 3D Simulator พร้อมระบบตรวจจับมือเปล่า (MediaPipe Hands) ช่วยให้ผู้เรียนสามารถทดลองปรับแต่งตัวแปรทางฟิสิกส์ด้วยมือเปล่าแบบไร้สัมผัส
    \item \textbf{การประมวลผลด้วยปัญญาประดิษฐ์เชิงลึกและภาษา Python (PINNs \& Qiskit):} นำเสนอการประยุกต์โครงข่ายประสาทเทียมตามหลักฟิสิกส์ (Physics-Informed Neural Networks) ในการหาผลเฉลยของสมการชเรอดิงเงอร์ และการจำลองวงจรควอนตัมด้วย Qiskit
\end{enumerate}

ผู้เขียนหวังเป็นอย่างยิ่งว่า ตำราเล่มนี้จะมีส่วนช่วยเปิดมิติใหม่ในการศึกษากลศาสตร์ควอนตัม ทำให้แนวคิดที่เป็นนามธรรมสูงกลายเป็นสิ่งที่จับต้องได้ สร้างแรงบันดาลใจ และเตรียมนักศึกษาให้มีความพร้อมต่อการขับเคลื่อนวิทยาศาสตร์และเทคโนโลยีขั้นสูงของประเทศต่อไป

\vspace{1.5cm}
\begin{flushright}
    ผู้ช่วยศาสตราจารย์ ดร.ชีวะ ทัศนา\\
    สาขาวิชาดิจิทัลการเกษตรและสิ่งแวดล้อม\\
    คณะวิทยาศาสตร์และเทคโนโลยี มหาวิทยาลัยราชภัฏรำไพพรรณี
\end{flushright}
